\documentclass[11pt,a4paper]{article}
\usepackage[T1]{fontenc}
\usepackage[utf8]{inputenc}
\usepackage[provide=*,french]{babel}
\usepackage{mathptmx}
\usepackage[margin=2.5cm]{geometry}
\usepackage{microtype,booktabs,longtable,array,enumitem,fancyhdr}
\usepackage[hidelinks]{hyperref}
\setlength{\parindent}{0pt}\setlength{\parskip}{7pt}
\setlength{\emergencystretch}{3em}
\setlist{nosep}
\pagestyle{fancy}\fancyhf{}\fancyhead[L]{\small Master Finance Verte -- Module 5}\fancyhead[R]{\small Évaluation finale}\fancyfoot[C]{\thepage}
\setlength{\headheight}{14pt}
\title{\Large Gouvernance d’entreprise, conformité et finance durable\\[8pt]\LARGE Audit critique du dossier SolarOne Maroc}
\author{Badr ABOUFARASSE, PhD}
\date{Version 1.0 -- 6 octobre 2026}
\begin{document}\maketitle
\begin{abstract}Évaluation intégratrice en binôme, à distance et sans oral. L’examen porte sur un dossier fictif de financement vert : analyse des preuves, hiérarchisation des constats et plan de correction. Le dossier ne constitue ni une consultation juridique ni une appréciation d’une entreprise réelle.\end{abstract}
\tableofcontents\newpage

\section{Votre mission}
Vous intervenez comme équipe d’examen préalable pour Atlas Transition, investisseur institutionnel fictif. Vous examinez le dossier SolarOne Maroc avant un engagement de 300 millions de MAD dans une société de projet solaire dont l’investissement total prévu est de 2 milliards de MAD. Le comité attend un audit critique et un plan de correction hiérarchisé.

Le dossier financier a déjà été analysé : dans les hypothèses retenues, le projet satisfait les critères financiers internes. Ces résultats ne constituent pas une validation de sa gouvernance, de sa conformité ou de ses caractéristiques environnementales. Aucun calcul financier supplémentaire n’est demandé.

Votre conclusion doit préciser si le dossier peut être accepté en l’état, accepté sous conditions, reporté pour instruction complémentaire ou rejeté. Plusieurs conclusions sont recevables si elles sont cohérentes avec les pièces, les incertitudes et les conditions proposées.

Toutes les organisations, personnes, chiffres, règles internes et citations de pièces du cas sont inventés pour l’évaluation. Les dates J-10, J-5 et J0 décrivent le calendrier fictif de la décision et ne sont pas des dates de remise des travaux. Les nouvelles pièces constituent une photographie propre à cette évaluation ; elles priment sur les variantes pédagogiques utilisées pendant les séances.

\section{Modalités et rendu}
Travail en binôme, entièrement à distance, sans soutenance orale. Documents du cours et sources institutionnelles autorisés. Charge de travail indicative : cinq à six heures par étudiant. Date limite de remise : à annoncer par l’enseignant ; la date du module précédent ne s’applique pas.

Un seul dépôt par binôme : rapport de quatre à cinq pages hors page de garde, références et annexes. Annexes : tableau des constats (deux à trois pages), plan de correction et de suivi (une à deux pages), puis une demi-page de contribution personnelle par étudiant. Le tableau peut être présenté en paysage.

La page de garde indique le module, le titre du cas, les noms des deux étudiants et la date. Utilisez une police lisible de 11 ou 12 points et numérotez les pages. Nom conseillé : M5\_Audit\_Nom1\_Nom2.pdf. Le PDF est recommandé ; Word ou une archive contenant le rapport et ses annexes sont également acceptés. Un fichier de 50 Mo au maximum.

L’utilisation éventuelle d’un outil d’IA doit être déclarée en fin de rapport : outil, usages et vérifications réalisées. Chaque référence et chaque argument restent sous la responsabilité des auteurs. Une réponse générique, des références inexistantes ou des faits ajoutés sans fondement ne permettent pas de satisfaire les critères de l’évaluation.

\section{Travail demandé}
1. Définir le périmètre de l’examen. Distinguez les décisions de la société de projet des décisions de l’investisseur. Identifiez l’organe compétent, les délégations, les fonctions de préparation, de surveillance et d’assurance indépendante.

2. Établir au moins huit constats distincts et non redondants couvrant les six thèmes du module. Pour chaque constat, citez une pièce et un passage ou une ligne ; indiquez le principe mobilisé, le statut de la preuve, le risque, l’information complémentaire et la correction proposée. Ne multipliez pas artificiellement les constats en séparant les conséquences d’un même problème.

3. Hiérarchiser les constats. Retenez trois niveaux : critique, majeur, modéré. Justifiez chaque priorité selon la gravité, l’urgence, la capacité de correction et les conséquences possibles pour la décision. Une information absente ne suffit pas à affirmer qu’une infraction ou une fraude est établie.

4. Développer les trois constats prioritaires dans le rapport. Expliquez leurs causes et leurs interactions : par exemple, comment une responsabilité mal répartie peut compromettre le contrôle puis la communication des résultats. Comparez une correction immédiate à une correction structurelle.

5. Examiner le statut des référentiels. Distinguez droit applicable, engagement contractuel de l’investisseur, politique interne et référence volontaire. Utilisez l’hypothèse de périmètre de P5 ; n’affirmez pas qu’un texte européen ou une norme internationale s’impose automatiquement à une société marocaine.

6. Proposer un plan opérationnel. Pour chaque mesure, précisez le responsable, l’échéance relative au financement, la preuve de clôture et l’organe qui vérifie cette clôture. Séparez les conditions préalables au premier versement, les engagements suivis après investissement et les motifs d’escalade.

7. Formuler une conclusion proportionnée. Précisez ce qui est démontré, ce qui reste incertain et quelle nouvelle information ferait changer votre décision. Proposez un dispositif d’engagement actionnarial et une réponse en cas de non-respect des engagements.

\section{Structure recommandée du rapport}
Résumé exécutif (une demi-page) : opinion sur l’état du dossier, trois priorités et recommandation. Périmètre et méthode (une demi-page) : corpus, statut des règles et limites. Analyse argumentée (deux à trois pages) : constats prioritaires, interactions et éléments satisfaisants à préserver. Conclusion et suivi (une page) : conditions, responsabilités et escalade.

L’annexe A est le tableau de tous les constats. L’annexe B rassemble le plan de correction et les indicateurs de suivi. L’annexe C contient la contribution de chaque étudiant : une analyse réellement menée, une difficulté rencontrée et une information qui pourrait modifier son jugement. Le rapport doit rester compréhensible sans recopier tout le dossier.

\section{Méthode et niveau d’exigence}
Statuts de preuve à utiliser : fait documenté ; contradiction entre pièces ; information manquante ; hypothèse à vérifier. Présentez la source de votre interprétation et signalez ses limites. Un constat doit relier un critère, une situation observée et une conséquence possible.

Le risque inhérent décrit l’exposition avant les mesures de maîtrise ; le risque résiduel décrit l’exposition après des mesures effectivement mises en œuvre. Un contrôle annoncé, une déclaration de conformité ou une mesure simplement prévue ne démontre pas son efficacité.

Une correction exploitable associe une action précise, un responsable, un délai et une preuve vérifiable. « Renforcer la transparence » ne suffit pas : indiquez quelle information sera produite, revue et communiquée. N’inventez ni autorisation administrative, ni seuil technique taxonomique, ni certification absents du dossier.

Un audit critique examine aussi les éléments satisfaisants. Expliquez quels dispositifs doivent être conservés et sous quelles conditions ils sont crédibles. La qualité est appréciée sur la solidité du raisonnement et la traçabilité, non sur le nombre maximal d’anomalies détectées.

\section{Les douze pièces du dossier}
Les pièces et les citations qui suivent sont intégralement fictives.

\subsection{P1 -- Note d’investissement et calendrier}
\textit{Direction des investissements d’Atlas Transition · J-10}

\textbf{P1.1.} Engagement envisagé : 300 millions de MAD. Participation minoritaire assortie d’un siège au conseil et d’un droit contractuel d’information trimestrielle. Première tranche prévue dix jours après la décision du comité. La société de projet gère la construction et l’exploitation ; Atlas Transition décide de l’engagement de ses propres capitaux.

\textbf{P1.2.} Le financement total du projet est de 2 milliards de MAD. La note financière conclut à l’admissibilité selon les critères internes, sous les hypothèses du modèle fourni pendant l’instruction. Le directeur du projet demande une décision rapide pour préserver le calendrier de construction.

\textbf{P1.3.} Le protocole d’investissement prévoit une vérification des conditions préalables avant décaissement. L’annexe précisant ces conditions est encore en préparation. Le conseil de SolarOne a approuvé le budget technique ; cette approbation est jointe à la note comme « preuve de validation globale du financement ».

\subsection{P2 -- Délégations et organisation}
\textit{Secrétariat général d’Atlas Transition · document en vigueur}

\textbf{P2.1.} Le comité d’investissement peut autoriser un engagement jusqu’à 250 millions de MAD. Au-delà de ce montant, il émet une recommandation au conseil d’administration de l’investisseur, qui autorise l’engagement. Aucun fractionnement en tranches ne modifie ce seuil, apprécié sur l’engagement total.

\textbf{P2.2.} La direction des investissements prépare le dossier. La fonction risques donne un avis distinct ; la conformité vérifie les partenaires et les engagements. Les deux fonctions peuvent saisir directement le président du comité. L’audit interne rend compte au comité d’audit de l’investisseur et examine la fiabilité des processus sans préparer les données opérationnelles.

\textbf{P2.3.} L’organigramme opérationnel joint au dossier place toutefois la responsable conformité sous l’autorité exclusive du directeur du projet. Une note manuscrite lui attribue aussi la préparation des indicateurs ESG et leur validation finale. La note ne mentionne aucun suppléant ni mécanisme de revue distincte.

\subsection{P3 -- Procès-verbal et déclaration d’intérêts}
\textit{Réunion du comité d’Atlas Transition · J-5}

\textbf{P3.1.} Le procès-verbal indique : « Le comité approuve définitivement l’engagement de 300 millions de MAD. La validation du conseil sera sollicitée au trimestre suivant. Le président est autorisé à signer sans autre condition. » Aucun avis signé de la fonction risques n’est annexé.

\textbf{P3.2.} Un membre du comité déclare détenir 20 \% d’EcoServ, entreprise proposée pour la maintenance de SolarOne. Il a participé à la présentation de l’offre et au vote. Le procès-verbal mentionne la déclaration, mais ni retrait de la délibération ni revue de l’offre par un tiers.

\textbf{P3.3.} EcoServ a présenté un prix compétitif et une expérience documentée sur deux installations comparables. Les deux autres offres n’ont pas été conservées dans le dossier transmis. La règle interne prévoit une déclaration d’intérêts, une décision sur le traitement du conflit et une trace des comparaisons pertinentes.

\subsection{P4 -- Registre des risques et fiche de contrôle}
\textit{Direction du projet · J-4}

\textbf{P4.1.} Trois lignes du registre : conflit d’intérêts, risque résiduel faible, mesure « déclaration du membre » ; qualité des données, risque résiduel faible, mesure « responsable ESG désignée » ; tension sur la ressource en eau, risque résiduel modéré, mesure « étude à réaliser après financement ». Les cases preuve d’efficacité et date de revue sont vides.

\textbf{P4.2.} Le comptable prépare chaque mois le rapprochement des paiements et un responsable financier distinct le signe. Les contrôles de mars à juin sont archivés et aucune exception ouverte n’est signalée. La fiche du même mois ne prévoit pas de contrôle équivalent pour le tableau ESG.

\textbf{P4.3.} Un courriel demande à l’audit interne de reconstituer les données manquantes de mai, de valider le tableau et de préparer ensuite un rapport indépendant sur ce même tableau.

\subsection{P5 -- Fiche des référentiels et de leur périmètre}
\textit{Service juridique · hypothèses expressément fixées pour le cas}

\textbf{P5.1.} SolarOne est une société de projet marocaine non cotée. Pour cet exercice, aucune obligation directe d’appliquer les règles européennes de publication d’informations de durabilité n’est établie. L’identification de ses obligations locales reste un point à documenter ; le dossier ne contient pas une consultation juridique complète.

\textbf{P5.2.} Atlas Transition exige contractuellement un examen de la qualification environnementale inspiré de la taxonomie européenne ainsi qu’une information structurée autour de la gouvernance, de la stratégie, des risques et des indicateurs. Ces exigences sont des conditions de l’investisseur dans le cas. Elles ne valent pas certification ni conclusion sur le champ juridique des textes.

\textbf{P5.3.} Le fournisseur de données a inscrit dans le formulaire : « L’énergie solaire étant verte, SolarOne est automatiquement conforme à la CSRD, à IFRS S2 et à la taxonomie européenne. » Aucun document de détermination du périmètre n’est joint.

\subsection{P6 -- Fiche de qualification environnementale}
\textit{Prestataire de données · version provisoire}

\textbf{P6.1.} Activité décrite : production d’électricité solaire. Objectif revendiqué : atténuation du changement climatique. Les champs figurent sous quatre rubriques : contribution substantielle ; absence de préjudice important aux autres objectifs ; garanties minimales ; critères techniques applicables.

\textbf{P6.2.} Le tableau porte « validé » pour toutes les rubriques. La contribution est justifiée par le type d’énergie. Les rubriques eau et biodiversité renvoient à une étude future. Les garanties minimales renvoient à une déclaration du directeur. Les critères techniques et leur version ne sont pas renseignés.

\textbf{P6.3.} Le taux annoncé est « 100 \% aligné ». Le dossier ne précise ni l’indicateur auquel ce taux se rapporte, ni le périmètre, ni le dénominateur, ni la personne ayant revu la fiche. Votre mission ne consiste pas à calculer un taux alternatif : formulez une conclusion sur ce que les pièces permettent effectivement de soutenir.

\subsection{P7 -- Tableau ESG et traçabilité}
\textit{Responsable ESG · données mensuelles déjà rapprochées pour l’exercice}

\textbf{P7.1.} Eau de nettoyage : avril, compteur 850 m³ et achats livrés 150 m³ ; mai, compteur non relevé et achats livrés 160 m³ ; juin, compteur 900 m³ et achats livrés 180 m³. Le total mensuel affiché en juin est 900 m³. Le prestataire confirme que l’eau livrée est distincte de l’eau mesurée au compteur : il n’y a pas de double compte dans ces données.

\textbf{P7.2.} Le total exact de juin fourni par l’équipe de rapprochement est 1 080 m³, au-dessus de la cible interne fictive de 1 000 m³. Le tableau public est resté à 900 m³ et porte « cible respectée ». Aucune estimation acceptée du compteur de mai n’est disponible ; cette valeur ne peut pas être supposée nulle.

\textbf{P7.3.} Le classeur est envoyé par courriel à plusieurs destinataires. Sa page de garde porte « version finale », alors que l’historique n’identifie ni la version diffusée ni une revue indépendante. Les sources de juin sont archivées ; celles du compteur de mai n’ont pas été produites.

\subsection{P8 -- Note climat et assurance}
\textit{Prestataire climat et mission de revue}

\textbf{P8.1.} La note annonce 80 000 tonnes de CO2e évitées par an par comparaison avec une référence électrique. Le scénario de référence, la période et les hypothèses de production ne sont pas joints. Un inventaire partiel mentionne 1 200 tonnes de CO2e pour les scopes 1 et 2 ; les équipements et la chaîne d’approvisionnement ne sont pas couverts dans la note.

\textbf{P8.2.} La communication soustrait les émissions évitées de l’inventaire et conclut à une « empreinte carbone négative certifiée ». Il n’existe dans le dossier ni conclusion sur la neutralité du projet ni certificat correspondant à cette expression.

\textbf{P8.3.} Une mission de revue externe signée existe. Son périmètre est limité aux factures d’électricité et aux relevés de juin ; elle ne couvre ni le mois de mai, ni les émissions évitées, ni l’alignement taxonomique, ni les garanties sociales. Sa conclusion précise ces exclusions.

\subsection{P9 -- Partenaire et prestation d’intermédiation}
\textit{Conformité · dossier partenaire}

\textbf{P9.1.} Une société dénommée Horizon Conseil demande une rémunération au succès de 2 \% du financement pour faciliter des démarches locales. Le contrat détaille le calendrier mais pas les prestations vérifiables ; la facture est payable sur un compte appartenant à une autre société.

\textbf{P9.2.} Le dossier contient un extrait d’immatriculation et une déclaration d’intégrité. L’identité des bénéficiaires effectifs n’est pas documentée ; la vérification des liens d’intérêts et des éventuels signaux de risque n’est pas achevée. Un courriel propose de verser l’avance puis de compléter les vérifications.

\textbf{P9.3.} Aucun élément du dossier ne démontre un paiement illicite ou une sanction. Il vous appartient de définir les diligences proportionnées et les conditions de paiement sans transformer un signal de risque en accusation établie.

\subsection{P10 -- Étude environnementale et réclamations}
\textit{Extraits du dossier opérationnel}

\textbf{P10.1.} L’étude initiale décrit le site et les mesures de gestion des déchets de chantier. Un budget de suivi est identifié. Une étude complémentaire sur le nettoyage, les prélèvements d’eau en période sèche et les habitats du site doit encore être réalisée. L’absence de ces compléments ne démontre pas à elle seule un dommage réalisé.

\textbf{P10.2.} Le registre comporte deux réclamations concernant l’accès à une piste et une réclamation d’un sous-traitant sur les conditions de travail. Deux accusés de réception sont archivés ; aucune réponse de clôture ni mécanisme d’escalade n’est indiqué. La fiche du prestataire conclut néanmoins « aucune controverse et garanties sociales validées ».

\textbf{P10.3.} Les autorisations requises doivent être inventoriées et vérifiées avant le premier versement selon le protocole interne. Le dossier transmis ne contient pas cet inventaire. N’inventez pas la nature exacte d’une autorisation locale : demandez son identification et sa justification juridique.

\subsection{P11 -- Projet de communication}
\textit{Direction de la communication · publication prévue à J+2}

\textbf{P11.1.} Texte proposé : « SolarOne : un investissement 100 \% durable, sans impact négatif, certifié neutre en carbone et intégralement aligné sur les standards européens. Toutes les données environnementales sont auditées. » Le texte ne précise aucun périmètre ni limite.

\textbf{P11.2.} Le responsable ESG recommande de préciser que les données sont provisoires et que la revue externe ne couvre qu’une partie du dossier. Le directeur du projet demande de conserver la version courte pour le lancement commercial. Aucune procédure formalisée de validation des allégations n’est jointe.

\textbf{P11.3.} Proposez dans votre rapport une reformulation courte fondée uniquement sur les faits documentés. Vous pouvez suspendre une allégation si les preuves ne permettent pas de la soutenir.

\subsection{P12 -- Projet d’engagement actionnarial}
\textit{Atlas Transition · annexe de suivi}

\textbf{P12.1.} Le projet de lettre prévoit des échanges trimestriels avec SolarOne et une amélioration des données ESG dans l’année. Les engagements ne comportent ni livrable précis, ni échéance individuelle, ni indicateur de clôture. Aucun calendrier n’est attaché à l’étude complémentaire sur l’eau.

\textbf{P12.2.} Le droit d’information et le siège au conseil sont prévus au protocole. Une escalade peut être proposée à l’organe compétent de l’investisseur, mais aucune séquence de dialogue, revue des engagements, décision sur une tranche non versée ou exercice des droits n’a été rédigée.

\textbf{P12.3.} Après le premier versement, l’investisseur ne dispose pas d’un droit unilatéral général de reprendre les fonds. Toute action doit respecter les droits contractuels. Proposez un suivi crédible en distinguant engagement, surveillance, escalade et décision sur les versements futurs.

\clearpage\section{Grille de notation}
\begin{longtable}{p{.30\textwidth}p{.08\textwidth}p{.53\textwidth}}\toprule Critère & Points & Attentes\\\midrule\endhead
Gouvernance et responsabilités & 4 & P1--P3 : compétence des organes, délégation, traitement des intérêts, séparation des rôles.\\[5pt]
Risques, contrôle interne et conformité & 4 & P4, P9--P10 : preuve des contrôles, risques résiduels, diligences proportionnées et priorités.\\[5pt]
Périmètre des référentiels et qualification & 4 & P5--P6 : droit, contrat et référence volontaire ; qualification, alignement et limites des preuves.\\[5pt]
Information, allégations et assurance & 3 & P7--P8, P11 : traçabilité, périmètre de revue et communication fidèle aux données.\\[5pt]
Correction, engagement et conclusion & 3 & P12 et ensemble du dossier : conditions, responsables, échéances, preuves de clôture et escalade.\\[5pt]
Qualité du rapport et contribution personnelle & 2 & Clarté, citations des pièces, références vérifiables, limites et contribution de chaque membre.\\[5pt]
\bottomrule\end{longtable}
Total : 20 points. Maîtrise élevée : analyse précise, fondée et opérationnelle ; satisfaisante : correcte mais partielle ; fragile : descriptive ; insuffisante : confusions majeures ou absence de preuves. Une conclusion alternative solidement justifiée est recevable.

\section{Modèles des annexes}
Annexe A : identifiant du constat ; pièce et passage ; principe ; statut de preuve ; risque ; priorité justifiée ; information complémentaire ; correction.

Annexe B : constat ; action ; responsable ; échéance ; preuve de clôture ; réviseur ; indicateur de suivi ; escalade.

Annexe C : pour chaque étudiant, une demi-page présentant sa contribution et une information susceptible de modifier son jugement. Un modèle Word modifiable est fourni sur la plateforme.

\section{Dépôt et confirmation}
Un seul membre dépose un fichier pour le binôme. Formats acceptés : PDF, Word, ZIP ou RAR, dans la limite de 50 Mo. Indiquer les noms, une adresse e-mail valide et le fichier ; attendre la confirmation et conserver le reçu. Date limite : à annoncer par l’enseignant.

\href{https://aboufarasse.github.io/Finance-Verte-Gouvernance-Conformite/evaluation-finale/}{Accéder à l’espace d’évaluation et de dépôt du module 5.}

\section{Références autorisées}
La présence d’une référence ne préjuge pas de son applicabilité juridique au cas. Utiliser la fiche P5 pour définir le périmètre.
\begin{enumerate}
\item OCDE (2023), Principes de gouvernance d’entreprise du G20 et de l’OCDE. Chapitres sur le conseil, la transparence et la durabilité. \href{https://doi.org/10.1787/734e2284-fr}{Consulter la source institutionnelle.}
\item The IIA (2026), Statements of Position : Three Lines Model. Rôles de gestion, surveillance, assurance et conseil. \href{https://www.theiia.org/en/resources/statements-of-position}{Consulter la source institutionnelle.}
\item IFRS Foundation, IFRS S1 : organisation de l’information sur les risques et opportunités de durabilité. \href{https://www.ifrs.org/issued-standards/ifrs-sustainability-standards-navigator/ifrs-s1-general-requirements/}{Consulter la source institutionnelle.}
\item IFRS Foundation, IFRS S2 : informations relatives au climat. \href{https://www.ifrs.org/supporting-implementation/supporting-materials-for-ifrs-sustainability-disclosure-standards/ifrs-s2/}{Consulter la source institutionnelle.}
\item Commission européenne, EU taxonomy for sustainable activities : objectifs, critères et actes applicables. \href{https://finance.ec.europa.eu/sustainable-finance/tools-and-standards/eu-taxonomy-sustainable-activities_en}{Consulter la source institutionnelle.}
\item GHG Protocol, Corporate Standard : inventaire des émissions de l’organisation. \href{https://ghgprotocol.org/corporate-standard}{Consulter la source institutionnelle.}
\item GHG Protocol (2019), Estimating and Reporting Avoided Emissions : comparaison avec une référence. \href{https://ghgprotocol.org/estimating-and-reporting-avoided-emissions}{Consulter la source institutionnelle.}
\end{enumerate}
\end{document}
